\documentclass{article}
\usepackage{graphicx} % Required for inserting images
\usepackage{amsmath}
\usepackage{setspace}
\onehalfspacing
\usepackage{tikz}
\usepackage{pgfplots}
\pgfplotsset{compat=1.18}

\title{CS556 - Derivatives Homework}
\author{John Rizzo}
\date{October 9, 2025}

\usepackage[margin=2cm]{geometry}

\begin{document}

\maketitle
\noindent
\vspace{1cm}
\textbf{Question 1} (20 points): Use the limit definition of the derivative to exactly evaluate the derivative:

\begin{enumerate}
    \item f(x) = $\sqrt{x+4}$
    \item f(x) = $\frac{3}{x}$
\end{enumerate}

\textbf{Solution:}
\begin{enumerate}
    \item For $f(x) = \sqrt{x+4}$, the derivative using the limit definition is:
    \[
    f'(x) = \lim_{h \to 0} \frac{\sqrt{(x+h)+4} - \sqrt{x+4}}{h}
    \]
    Rationalizing the numerator:
    \[
    = \lim_{h \to 0} \frac{(\sqrt{x+h+4} - \sqrt{x+4})(\sqrt{x+h+4} + \sqrt{x+4})}{h(\sqrt{x+h+4} + \sqrt{x+4})}
    \]
    \[
    = \lim_{h \to 0} \frac{(x+h+4) - (x+4)}{h(\sqrt{x+h+4} + \sqrt{x+4})}
    \]
    \[
    = \lim_{h \to 0} \frac{h}{h(\sqrt{x+h+4} + \sqrt{x+4})}
    \]
    \[
    = \lim_{h \to 0} \frac{1}{\sqrt{x+h+4} + \sqrt{x+4}} = \frac{1}{2\sqrt{x+4}}
    \]

    Thus, $f'(x) = \frac{1}{2\sqrt{x+4}}$.

    \item For $f(x) = \frac{3}{x}$, the derivative using the limit definition is:
    \[
    f'(x) = \lim_{h \to 0} \frac{\frac{3}{x+h} - \frac{3}{x}}{h}
    \]
    Finding a common denominator:
    \[
    = \lim_{h \to 0} \frac{\frac{3x - 3(x+h)}{(x+h)x}}{h}
    \]
    \[
    = \lim_{h \to 0} \frac{\frac{-3h}{(x+h)x}}{h}
    \]
    Simplifying:
    \[
    = \lim_{h \to 0} \frac{-3}{(x+h)x} = -\frac{3}{x^2}
    \]

    Thus, $f'(x) = -\frac{3}{x^2}$.
\end{enumerate}

\vspace{1cm}
\textbf{Question 2} (\textit{20 points}) Find the derivative of the following functions:

\begin{enumerate}
    \item $f(x) = 3x^3 - \frac{4}{x^2} = 9x^2 - 4x^{-2}$
    \item $f(x) = (4 - x^2)^3$ = $u^3$ where $u = 4 - x^2$, so $f'(u) = 3u^2$ and $u' = -2x$ so by the chain rule $f'(x) = f'(u) \cdot u' = 3(4 - x^2)^2 \cdot (-2x) = -6x(4 - x^2)^2$
    \item $f(x) = e^{sin(x)}$
    $f'(x) = e^{sin(x)} \cdot cos(x)$ by the chain rule.
    \item $f(x) = ln(x+2)$
    $f'(x) = \frac{1}{x+2}$
    \item $f(x) = x^2cos(x) + xtan(x)$ \\
    Using the product rule:
    $f'(x) = 2xcos(x) - x^2sin(x) + tan(x) + xsec^2(x)$
    \item $f(x) = \sqrt{3x^2+2}$ \\
    $f'(x) = \frac{1}{2\sqrt{3x^2+2}} \cdot 6x = \frac{6x}{2\sqrt{3x^2+2}} = \frac{3x}{\sqrt{3x^2+2}}$
    \item $f(x) = \frac{x}{4}\sin^{-1}(x)$ \\
    Using the product rule:
    $f'(x) = \frac{1}{4}\sin^{-1}(x) + \frac{x}{4} \cdot \frac{1}{\sqrt{1 - x^2}} = \frac{\sin^{-1}(x)}{4} + \frac{x}{4\sqrt{1 - x^2}}$

    \item $x^2y = (y+2) + xy\sin(x)$ \\
    $y = \frac{(y+2) + xy\sin(x)}{x^2}$ \\
    Calculate the derivative of both sides \\
    Left Side: $\frac{d}{dx}(x^2y) = \frac{d}{dx}(x^2) \cdot y + x^2 \cdot \frac{dy}{dx} = 2xy + x^2 \frac{dy}{dx}$ \\
    Right Side: $\frac{d}{dx}((y+2) + xy\sin(x)) = \frac{dy}{dx} + y\sin(x) + x\frac{dy}{dx}\sin(x) + xy\cos(x)$ \\
    Setting the derivatives equal:
    $2xy + x^2 \frac{dy}{dx} = \frac{dy}{dx} + y\sin(x) + x\frac{dy}{dx}\sin(x) + xy\cos(x)$ \\
    Rearranging to isolate $\frac{dy}{dx}$:
    $x^2 \frac{dy}{dx} - \frac{dy}{dx} - x\frac{dy}{dx}\sin(x) = y\sin(x) + xy\cos(x) - 2xy$ \\
    $\frac{dy}{dx}(x^2 - 1 - x\sin(x)) = y\sin(x) + xy\cos(x) - 2xy$ \\
    Finally:
    $\frac{dy}{dx} = \frac{y\sin(x) + xy\cos(x) - 2xy}{x^2 - 1 - x\sin(x)}$
\end{enumerate}

\vspace{1cm}
\textbf{Question 3} (\textit{20 points}) The following questions consider the wind speeds of Hurricane Katrina, which affected New Orleans, Lousiana, in August 2005.  The data are displayed in a table.

\begin{table}[htbp]
    \centering
    \caption{Wind Speeds of Hurricane Katrina on August 26, 2005}

    \begin{tabular}{|c|c|}
        \hline
        Hours After Midnight & Wind Speed (mph) \\
        \hline
        1 & 45 \\
        \hline
        5 & 75 \\
        \hline
        11 & 100 \\
        \hline
        29 & 115 \\
        \hline
        49 & 145 \\
        \hline
        58 & 175 \\
        \hline
        73 & 155 \\
        \hline
        81 & 125 \\
        \hline
        85 & 95 \\
        \hline
        107 & 35 \\
        \hline
    \end{tabular}
\end{table}

\begin{enumerate}
    \item Using the table, estimate the deriative of the wind speed at hour 39.  What is the physical meaning?

    To estimate the derivative at hour 39, we can use the data points at hours 29 and 49:
    \[f'(39) \approx \frac{f(49) - f(29)}{49 - 29} = \frac{145 - 115}{20} = \frac{30}{20} = 1.5 \text{ mph per hour}\]
    The physical meaning of this derivative is that at hour 39, the wind speed of Hurricane Katrina was increasing at a rate of approximately 1.5 miles per hour for each hour.

    \item Estimate the derivative of the wind speed at hour 83.  What is the physical meaning?
    
    To estimate the derivative at hour 83, we can use the data points at hours 81 and 85:
    \[f'(83) \approx \frac{f(85) - f(81)}{85 - 81} = \frac{95 - 125}{4} = \frac{-30}{4} = -7.5 \text{ mph per hour}\]       
    The physical meaning of this derivative is that at hour 83, the wind speed of Hurricane Katrina was decreasing at a rate of approximately 7.5 miles per hour for each hour.
\end{enumerate}

\vspace{1cm}
\textbf{Question 4} (\textit{20 points}) The famous Regiomontanus' problem for angle maximization was proposed during the 15th centrury.  A painting hangs on a wall with the bottom of the painting a distance $a$ feet above eye level, and the top $b$ feet above eye level. What distance $x$ (in feet) from the wall should the viewer stand to maximize the angle subtended by the painting, $\theta$?

\textbf{Solution:}
$tan(\theta - \alpha) = \frac{tan(\theta) - tan(\alpha)}{1 + tan(\theta)tan(\alpha)}$ \\
$ = \frac{\frac{b}{x} - \frac{a}{x}}{1 + \frac{b}{x} \cdot \frac{a}{x}}$ \\
$ = \frac{\frac{b-a}{x}}{1 + \frac{ab}{x^2}}$ \\
$ = \frac{(b-a)x}{x^2 + ab}$ \\
To maximize $\theta$, we maximize $tan(\theta)$: \\
Let $f(x) = \frac{(b-a)x}{x^2 + ab}$ \\
Using the quotient rule: \\
$f'(x) = \frac{(b-a)(x^2 + ab) - (b-a)x(2x)}{(x^2 + ab)^2}$ \\
$= \frac{(b-a)(ab - x^2)}{(x^2 + ab)^2}$ \\
Setting $f'(x) = 0$ to find critical points: \\
$(b-a)(ab - x^2) = 0$ \\
Since $b \neq a$, we have: \\
$ab - x^2 = 0$ \\
$x^2 = ab$ \\
$x = \sqrt{ab}$ \\
Thus, the viewer should stand a distance of $\sqrt{ab}$ feet from the wall to maximize the angle subtended by the painting.

\vspace{1cm}
\textbf{Question 5} (\textit{20 points}) An airline sells tickets from Tokyo to Detroit for \$1200.  There are 500 seats available and a typical flight books 350 seats.  For every \$10 decrease in price, the airline observes an additional five seats sold.  What should the fare be to maximize profit? How many passengers would be onboard?

\textbf{Solution:}
Let $x$ be the number of \$10 decreases in price.  Then the price per ticket is $1200 - 10x$ and the number of passengers is $350 + 5x$.  The profit function is:
\[P(x) = (1200 - 10x)(350 + 5x) \]
Expanding:
\[P(x) = 420000 + 6000x - 3500x - 50x^2 \]
\[P(x) = -50x^2 + 2500x + 420000 \]
To find the maximum profit, we can take the derivative of the profit function with respect to \( x \) and set it to zero:
\[ \frac{dP}{dx} = -100x + 2500 \]
Setting the derivative equal to zero:
\[ -100x + 2500 = 0 \implies x = 25 \]
Substituting \( x = 25 \) back into the price and passenger equations:
\[ \text{Price} = 1200 - 10 \times 25 = 950 \text{ dollars} \]
\[ \text{Passengers} = 350 + 5 \times 25 = 475 \]
Therefore, the airline should set the fare at \$950 to maximize profit, with 475 passengers onboard.

\end{document}